\documentclass[11pt]{article}
\usepackage{../../../tex/lecturestyle}

\title{Sperner's Theorem}
\author{}
\date{}

\begin{document}
\thispagestyle{fancy}

\lecturetitleband{Set Operations \& Functions}{Sperner's Theorem}{5}

\section{Overview}
Among all subsets of $\{0, 1, \ldots, n-1\}$, how many can be chosen so
that no one is contained in another? \textbf{Sperner's theorem} answers
this exactly: the maximum is $\binom{n}{\lfloor n/2 \rfloor}$, achieved by
simply taking every subset of size $\lfloor n/2 \rfloor$. This lecture
gives two independent proofs --- one via a clean probabilistic inequality
(LYM), one via an explicit combinatorial decomposition --- along with the
exact characterization of when equality holds and a useful matching
corollary. The ideas here (chains, antichains) are foundational for order
theory generally, including results like Dilworth's theorem.

\section{Chains and Antichains}
Write $B_n$ for the \emph{Boolean lattice}: all subsets of
$\{0,\ldots,n-1\}$, ordered by $\subseteq$. Drawing an edge $S \to T$
whenever $S \subsetneq T$ and $|T|=|S|+1$, with same-size subsets drawn at
the same height, gives $B_n$'s \emph{Hasse diagram}.

\begin{definitionbox}
A set $C \subseteq B_n$ is a \emph{chain} if every two of its elements are
comparable ($S \subseteq T$ or $T \subseteq S$). A set $A \subseteq B_n$
is an \emph{antichain} if every two \emph{distinct} elements are
incomparable.
\end{definitionbox}

\begin{theorembox}[(Chains meet antichains at most once)]
For any chain $C$ and antichain $A$, $|C \cap A| \le 1$.

\emph{Proof.} Two distinct elements of $C \cap A$ would have to be both
comparable (from $C$) and incomparable (from $A$). \qed
\end{theorembox}

Let $B_{n,k}$ denote the size-$k$ subsets of $\{0,\ldots,n-1\}$. Since
same-size sets are never comparable, $B_{n,k}$ is itself an antichain, of
size $\binom{n}{k}$; taking $k = \lfloor n/2 \rfloor$ maximizes this.

\begin{theorembox}[(Sperner)]
Every antichain $A \subseteq B_n$ satisfies
$|A| \le \binom{n}{\lfloor n/2 \rfloor}$.
\end{theorembox}

\section{Proof 1: The LYM Inequality}
For a permutation $p = (p_0, \ldots, p_{n-1})$ of $\{0,\ldots,n-1\}$,
define the \emph{maximal chain}
$C_p = \{\emptyset, \{p_0\}, \{p_0,p_1\}, \ldots, \{p_0,\ldots,p_{n-1}\}\}$
--- one subset of each size $0,\ldots,n$.

\begin{theorembox}[(Lubell--Yamamoto--Meshalkin inequality)]
For any antichain $A \subseteq B_n$,
\[
  \sum_{S \in A} \binom{n}{|S|}^{-1} \;\le\; 1.
\]

\emph{Proof.} Pick $p$ uniformly at random among all $n!$ permutations,
and consider $\mathbb{E}[\,|C_p \cap A|\,]$. Since $|C_p \cap A| \le 1$
always (previous theorem), the expectation is at most $1$. On the other
hand, for fixed $S$ with $|S|=m$, $C_p$ contains a unique $m$-element set,
uniformly distributed among all $\binom{n}{m}$ of them as $p$ varies, so
$\Pr[S \in C_p] = \binom{n}{m}^{-1}$. Summing this probability over
$S \in A$ gives $\mathbb{E}[\,|C_p \cap A|\,] = \sum_{S \in A}
\binom{n}{|S|}^{-1}$, which is therefore $\le 1$. \qed
\end{theorembox}

\subsection{Deriving Sperner's theorem}
Since $\binom{n}{k}$ is maximized at $k = \lfloor n/2 \rfloor$, every term
satisfies $\binom{n}{|S|}^{-1} \ge \binom{n}{\lfloor n/2\rfloor}^{-1}$.
Substituting into LYM:
\[
  1 \;\ge\; \sum_{S \in A} \binom{n}{|S|}^{-1}
  \;\ge\; |A| \cdot \binom{n}{\lfloor n/2 \rfloor}^{-1},
\]
so $|A| \le \binom{n}{\lfloor n/2 \rfloor}$.

\subsection{When is equality achieved?}
Equality forces every inequality along the way to be tight: every $S \in
A$ must satisfy $\binom{n}{|S|} = \binom{n}{\lfloor n/2\rfloor}$, i.e.\
$|S| \in \{\lfloor n/2\rfloor, \lceil n/2\rceil\}$, and $|C_p \cap A|=1$
for \emph{every} permutation $p$.

\begin{theorembox}
If $|A| = \binom{n}{\lfloor n/2\rfloor}$, then $A = B_{n,n/2}$ when $n$ is
even, and $A = B_{n,(n-1)/2}$ or $A = B_{n,(n+1)/2}$ when $n$ is odd.

\emph{Proof sketch (odd case; even is immediate).} Write $n=2m+1$. Every
$S \in A$ has size $m$ or $m+1$; suppose some $S \in A$ has $|S|=m$, and
let $x \in S$, $y \notin S$, $T = S \setminus\{x\} \cup \{y\}$. Some
permutation's maximal chain contains both $T$ and $S \cup \{y\}$ (both
have the same two elements $\{x,y\}$ swapped relative to a common
$(m{+}1)$-set); since $A$ meets that chain exactly once, and
$S\cup\{y\} \notin A$ (it strictly contains $S \in A$), we get $T \in A$.
Repeating this swap argument shows $A$ contains \emph{every} $m$-element
set once it contains one, so (matching $|A|=\binom{n}{m}$) $A = B_{n,m}$
exactly. \qed
\end{theorembox}

\section{Proof 2: Symmetric Chain Decomposition}
\begin{definitionbox}
A chain $S_0 \subsetneq S_1 \subsetneq \cdots \subsetneq S_m$ in $B_n$ is
\emph{symmetric} if consecutive sets differ in size by exactly $1$, and
$|S_0| + |S_m| = n$ (it is centered around the middle rank).
\end{definitionbox}

\begin{theorembox}
$B_n$ can be partitioned into disjoint symmetric chains.

\emph{Proof (induction on $n$).} $B_0 = \{\emptyset\}$ is trivially one
chain. Given a symmetric chain decomposition of $B_{n-1}$, write
$S^+ = S \cup \{n-1\}$ for $S \subseteq \{0,\ldots,n-2\}$. Each chain
$S_0 \subsetneq \cdots \subsetneq S_m$ of $B_{n-1}$ splits into two chains
of $B_n$: $S_0 \subsetneq \cdots \subsetneq S_m \subsetneq S_m^+$, and (if
$m \ge 1$) $S_0^+ \subsetneq \cdots \subsetneq S_{m-1}^+$. Both are
symmetric in $B_n$: the first has sizes $|S_0|$ and $|S_m|+1$, summing to
$n$; the second has sizes $|S_0|+1$ and $|S_{m-1}|+1 = |S_m|$, also
summing to $n$. Every subset of $\{0,\ldots,n-1\}$ is uniquely $S$ or
$S^+$ for some $S \subseteq \{0,\ldots,n-2\}$, so applying this split to
every chain of the $B_{n-1}$ decomposition covers all of $B_n$. \qed
\end{theorembox}

\subsection{Deriving Sperner's theorem, again}
Fix $k = \lfloor n/2 \rfloor$. Any symmetric chain has min size $a$ and
max size $b=n-a$ with $a \le b$, so $a \le k \le b$: every symmetric chain
passes through rank $k$ \emph{exactly once}. Hence the decomposition has
exactly $\binom{n}{k}$ chains, and (chains meet antichains at most once)
any antichain has size at most $\binom{n}{k}$.

\subsection{A matching corollary}
The same decomposition gives, for free, a maximum matching between
consecutive levels.

\begin{theorembox}
For $0 \le k < n$, there is an injective map pairing each of
$\min\bigl(\binom{n}{k}, \binom{n}{k+1}\bigr)$ sets in $B_{n,k}$ with a
proper superset in $B_{n,k+1}$.

\emph{Proof.} Take every symmetric chain that passes through both rank $k$
and rank $k+1$, and pair its rank-$k$ and rank-$(k+1)$ members. If
$k < n/2$, every chain reaching rank $k$ also reaches rank $k+1$ (since
its span is symmetric about the middle), giving $\binom{n}{k}$ pairs;
symmetrically for $k \ge n/2$. \qed
\end{theorembox}

\begin{notebox}
This matching can also be proved first (by other means) and used to
\emph{derive} Sperner's theorem: chain together matched pairs across all
levels to build $\binom{n}{\lfloor n/2\rfloor}$ chains covering $B_n$,
recovering the same bound.
\end{notebox}

\begin{notebox}
Sperner's theorem generalizes from $B_n$ (a product of $n$ two-element
chains) to a product of chains of arbitrary lengths, where the maximum
antichain is the ``middle layer'' by coordinate sum; the symmetric-chain
argument adapts with the same overall shape, though the details are
substantially more involved.
\end{notebox}

\section{Reference Implementations}
Both examples below implement the symmetric chain decomposition (verified
to be a valid partition of $B_n$ into exactly $\binom{n}{\lfloor
n/2\rfloor}$ symmetric chains), the level-matching corollary, the LYM
inequality (checked on random maximal antichains, with exact rational
arithmetic), and an independent cross-check of the whole theorem via
Dilworth's theorem: the maximum antichain size equals the minimum number
of chains needed to cover the poset, computed here by bipartite matching.

\subsection{Rust}
\lstinputlisting[language=Rust]{example.rs}

\subsection{Python}
\lstinputlisting[language=Python]{example.py}

\section{Worked Example}
\begin{examplebox}
For $n=3$, the inductive construction produces the symmetric chains
$\{\emptyset,\{0\},\{0,1\},\{0,1,2\}\}$, $\{\{2\},\{0,2\}\}$, and
$\{\{1\},\{1,2\}\}$ --- three chains, matching
$\binom{3}{1}=3=\binom{3}{\lfloor 3/2\rfloor}$. Reading off rank $1$ and
rank $2$ from each chain gives a perfect matching between $B_{3,1}$ and
$B_{3,2}$ (both size $3$): $\{0\}\!\to\!\{0,1\}$, $\{2\}\!\to\!\{0,2\}$,
$\{1\}\!\to\!\{1,2\}$ --- exactly what the reference implementations
compute and verify for every $n$ up to $8$--$10$.
\end{examplebox}

\section{Further Reading}
\begin{itemize}[leftmargin=1.5em]
  \item AtCoder Algorithm Lectures index (source of this curriculum):
        \url{https://info.atcoder.jp/entry/algorithm_lectures/index}
  \item E. Sperner, ``Ein Satz über Untermengen einer endlichen Menge'',
        Mathematische Zeitschrift 27(1), 1928, pp.\ 544--548.
  \item Wikipedia, \emph{Lubell--Yamamoto--Meshalkin inequality}:
        \url{https://en.wikipedia.org/wiki/Lubell\%E2\%80\%93Yamamoto\%E2\%80\%93Meshalkin_inequality}
  \item Wikipedia, \emph{Dilworth's theorem}:
        \url{https://en.wikipedia.org/wiki/Dilworth's_theorem}
\end{itemize}

\end{document}
